% =============================================================================
%  «ТЕОРИЯ СИСТЕМ» — ствол серии «Архитектура Рассуждения»
%  Блок 2. ОБЪЕКТ: каноническая система.  Разделы 2.1–2.8 + хорошая форма (черновик прозы).
%  Стиль: «выводим и показываем» — формулы/имена в тексте кратко; в каждом разделе
%  сначала ВЫВОД/ПОКАЗ, затем РЕЗУЛЬТАТЫ; репо-якоря — в приложении, не в сносках.
%  Самокомпилируемо: pdflatex blok-2-obekt.tex (см. память latex-build-windows).
% =============================================================================
\documentclass[12pt,a4paper]{article}
\usepackage{cmap}
\usepackage[T2A]{fontenc}
\usepackage[utf8]{inputenc}
\usepackage[russian]{babel}
\usepackage{paratype}
\usepackage{amsmath,amssymb,amsthm}
\usepackage{listings}
\usepackage[expansion=false]{microtype}
\usepackage[margin=28mm]{geometry}

\newcommand{\rezultat}{\medskip\noindent\emph{Что отсюда следует.}\ }

\title{\textbf{Теория Систем}\\[2mm]\large Блок 2. Объект: каноническая система}
\author{}
\date{}

\begin{document}
\maketitle
\thispagestyle{empty}

\noindent\small\textit{Черновик: разделы 2.1--2.8 и хорошая форма. Продолжает Блок~1 (вход):
там цепь $A=\exists \to$ различение $\to$ законы $\to$ принципы $\to$ E/R/R дала объект; здесь
объект разбирается изнутри. Машинные якоря --- в Приложении в конце блока, не в сносках.}\normalsize

\bigskip

\noindent Цепь предыдущего блока дала объект --- триаду E/R/R. Но «система есть тройка
Elements\,/\,Roles\,/\,Rules» было бы лишь \emph{переописанием}. Этот блок показывает иное:
система есть триада, скреплённая ещё двумя вещами --- \emph{гейтом правил} и \emph{градацией по
уровням}; триада \emph{одна} систему не определяет, и это результат, а не дефект.

% =============================================================================
\section*{2.1\quad Каноническая система: триада в одной форме}
% =============================================================================

Канонический объект теории несёт все три аспекта в одной записи. Правило здесь --- это условие
на пару $\langle$носитель, отношение$\rangle$: для данных элементов и данной их соотнесённости
оно либо выполнено, либо нет. Простейшие примеры --- правило \emph{тривиальное} (выполнено
всегда) и правило-\emph{эквивалентность} (отношение рефлексивно, симметрично, транзитивно).

Решающая точка --- та, где правило \emph{применено} к паре (Elements, Roles) и \emph{выполнено}.
Назовём это \emph{функциональным фактом}: именно здесь правила «гейтят» --- пропускают одни
конфигурации носителя и ролей и отвергают другие. Этот факт --- семя двух дальнейших
результатов: ранговой асимметрии (§2.4) и закона композиции (Блок~3).

\rezultat
Система --- не набор из трёх независимых вещей, а триада, \textbf{скреплённая функциональным
фактом}: правило, выполненное на своих элементах и ролях. Без этого скрепа три аспекта остались
бы лежать порознь и системы бы не составляли.

% =============================================================================
\section*{2.2\quad Экстракция: разбор, а не сборка}
% =============================================================================

Из готовой системы триада \emph{вытягивается}: \emph{что} в ней (Elements), \emph{как} соотнесено
(Roles), \emph{почему} так (Rules). Эпистемически это \emph{разбор} наличного объекта, а не сборка
его из трёх заготовленных частей --- порядок познания обратен порядку порождения (§1.5).

Три аспекта при этом взаимно-исключающи: элемент не есть правило, роль не есть элемент. Каждая
грань занимает свою категорию и не подменяет другую.

\rezultat
Первичен сам объект; его три аспекта --- грани, \emph{полученные разбором}. Поэтому говорить
«система состоит из E, R, R» неточно: точнее --- «в системе \emph{различимы} E, R, R».

% =============================================================================
\section*{2.3\quad Теорема представления}
% =============================================================================

Что же достаёт триаду до \emph{полной} системы? Два скрепа: \emph{функциональный факт} (гейт
правил --- это L4) и \emph{градация по уровням} (это P1). Триада плюс эти два --- вот система.

Здесь --- ядро блока, и оно honest. \textbf{Триада одна систему не определяет.} Можно предъявить
две системы, которые делят \emph{всю} триаду --- те же элементы, те же роли, те же правила, --- и
всё же различаются градацией по уровням. Значит триада \emph{необходима}, но \emph{не достаточна}.

Это и есть поправка к наивному «система $=$ тройка»: такое равенство было бы переописанием.
Честная теорема представления показывает, \emph{чего тройке не хватает}, --- гейта и градации.

\rezultat
Представление системы: \textbf{триада $\oplus$ гейт правил $\oplus$ градация по уровням.} Триада ---
\emph{структурное ядро}, необходимое, но не полное. (Единственность здесь --- на уровне данных;
вопрос о равенстве доказательств в саму запись не входит и отдельной опоры не требует.)

% =============================================================================
\section*{2.4\quad Ранговая асимметрия правил}
% =============================================================================

Наблюдение «правила первичнее ролей, роли первичнее элементов» в этом блоке становится
\emph{теоремой}. Правила \emph{гейтят}; а пара (элементы, роли) правила \emph{не определяет}: есть
две системы с одинаковыми элементами и ролями, но разными правилами. Значит знание $E$ и $R$ не
даёт $R$ (Rules) --- ранг правил строго выше.

Это и есть доказательство \emph{генеративного порядка} Rules${}\to{}$Roles${}\to{}$Elements,
заявленного в §1.5: правило задаёт возможные роли, роли --- что может стать элементом.

\rezultat
Первичность правил --- не лозунг, а \textbf{теорема}: $(E,R)$ недоопределяют $R$. Честно: это не
переформулировка функционального факта, а \emph{структурная} асимметрия --- она заложена в самих
типах записи, и из неё доказаны следствия.

% =============================================================================
\section*{2.5\quad Неприводимость триады}
% =============================================================================

Три аспекта --- три \emph{независимые} степени свободы: ни один не сводится к двум другим. Правила
независимы (это §2.4); роли независимы от элементов (одни и те же элементы допускают разные
соотнесения); элементы независимы --- у носителя может быть вовсе иной \emph{тип}.

\rezultat
«Три грани, ни одну не опустить» --- формализовано точно: как \textbf{независимость трёх осей}.
Опустить любую --- не упростить систему, а перестать иметь систему.

% =============================================================================
\section*{2.6\quad Три уровня бытия и Status}
% =============================================================================

Существование расслаивается \emph{лестницей пресуппозиции}. Существовать (нечто \emph{есть} --- снова
единица, §1.2) $\to$ принадлежать (нечто \emph{входит} в систему по критерию) $\to$ функционировать
(нечто \emph{играет} роль). Каждая ступень предполагает предыдущую: нельзя играть роль, не входя;
нельзя входить, не существуя.

Отсюда уточнение, которое легко спутать. \emph{Status} (чем нечто \emph{стало}, пройдя правила) ---
не то же, что \emph{Role} (\emph{почему} оно здесь). Status \emph{выведен}: он есть применённое
правило, а не отдельный примитив рядом с ролью.

\rezultat
Бытие \emph{градуально по пресуппозиции}: существование $\to$ принадлежность $\to$ функционирование.
А «статус» --- \emph{производное} от правила, не новая сущность. (Прежняя пометка-заглушка о «трёх
уровнях бытия» заменена этой реальной лестницей.)

% =============================================================================
\section*{2.7\quad Актуализация: подъём по уровню}
% =============================================================================

Уровни связаны \emph{операцией}. Продукты уровня $L$ становятся элементами уровня $L+1$, и при
подъёме триада переносится целиком: носитель, роли и правила сохраняются, а валидность по уровню
поднимается вместе с ними. Это \emph{общий} ход, а не частный случай чисел или рациональных.

\rezultat
Подъём по уровню --- \textbf{операция над системами}, а не свойство отдельного примера. Это
превью вертикали: башня уровней (Блок~6) собирается именно из этого шага.

% =============================================================================
\section*{2.8\quad Мост двух «систем»}
% =============================================================================

В теории живут две формы «системы». Ранняя --- \emph{критериальная}: нечто собрано по признаку
принадлежности; она служит блокировке парадоксов (критерий прежде системы, P2). И каноническая ---
\emph{триада} E/R/R. Первая \emph{отображается} во вторую: её члены становятся элементами, равенство
--- ролями, признак --- правилом.

\rezultat
Каноничен \textbf{FunctionalSystem}; критериальная форма --- его частный, \emph{обеднённый} случай.
Честно: это \emph{вложение}, а не изоморфизм --- критериальная форма тоньше (без собственных ролей и
правил), а её образ в триаде богаче.

% =============================================================================
\section*{2.9\quad Хорошая форма системы}
% =============================================================================

У системы есть \emph{разрешимый} признак хорошей формы: каждый компонент занимает \emph{ровно одну}
категорию E/R/R \emph{и} в системе нет само-референции (это P1 из §1.4). Признак вычислим: по самой
записи видно, хорошо ли она сформирована.

Числа проходят; проходит и шахматная партия. А «множество всех множеств, не содержащих себя»
(Рассел), «это высказывание ложно» (Лжец), «гетерологическое» (Греллинг) --- \emph{не} проходят: они
\emph{дурно сформированы}. Парадоксы здесь не «растворяются» --- они с самого начала \emph{невозможны}
как хорошо сформированные системы.

\rezultat
Хорошая форма --- \textbf{условие применимости} представления (§2.3): теорема о триаде говорит о
хорошо сформированных системах. Парадокс --- это \emph{нарушение} хорошей формы, а не загадка
внутри неё.

\medskip
\noindent\emph{Переход.} Объект разобран изнутри: триада, её скрепы, ранг, неприводимость, уровни
бытия, подъём, мост форм, хорошая форма. Следующий блок показывает, как такие объекты
\emph{складываются} --- морфизмы, произведения, факторы, категория систем.

% =============================================================================
\appendix
\section*{Приложение к Блоку 2 — машинные якоря}
% =============================================================================
\small\raggedright\sloppy
\noindent Всё перечисленное машинно проверено в репозитории; единственная аксиома ветви ---
\texttt{classic} (= Закон Исключённого Третьего, L3). Строки сверяются при финальной сборке.

\begin{itemize}\itemsep2pt
\item \textbf{§2.1--2.2.} Запись \texttt{Constitution} (правило $=\,\forall\,D\,R,\ \mathrm{Prop}$),
\texttt{FunctionalSystem} (6 полей, \emph{функциональный факт} \texttt{fs\_functional}); экстракторы
\texttt{get\_Elements/Roles/Rules}; взаимоисключение категорий \texttt{categories\_exclusive},
\texttt{err\_element\_not\_rule}.\\
\emph{Файлы:} \texttt{TheoryOfSystems\_Core\_ERR.v}~§XVII, \texttt{Roles.v}.

\item \textbf{§2.3 (флагман).} \texttt{triple\_not\_complete} (триада не определяет систему);
\texttt{fs\_eta}; капстоун \texttt{err\_representation}.\\
\emph{Файл:} \texttt{ERRRepresentation.v}.

\item \textbf{§2.4.} \texttt{rules\_are\_the\_gate}, \texttt{elements\_underdetermine\_roles},
\texttt{different\_rules}; капстоун \texttt{err\_rank\_asymmetry}.\\
\emph{Файл:} \texttt{ERRRankAsymmetry.v}.

\item \textbf{§2.5.} \texttt{rules\_independent}, \texttt{roles\_independent},
\texttt{elements\_independent}; капстоун \texttt{err\_triad\_irreducible}.\\
\emph{Файл:} \texttt{ERRIrreducible.v}.

\item \textbf{§2.6.} \texttt{three\_levels\_chain} (лестница пресуппозиции),
\texttt{status\_determined}, \texttt{status\_needs\_rule}.\\
\emph{Файл:} \texttt{ERRBeingAndStatus.v}.

\item \textbf{§2.7.} \texttt{fs\_lift}, \texttt{lift\_elements\_valid\_up}; капстоун
\texttt{err\_actualization}.\\
\emph{Файл:} \texttt{ERRActualization.v}.

\item \textbf{§2.8.} \texttt{system\_to\_FS} (вложение, не изо); капстоун
\texttt{err\_system\_bridge}.\\
\emph{Файл:} \texttt{ERRSystemBridge.v}.

\item \textbf{§2.9.} \texttt{well\_formedness\_decidable}, \texttt{russell\_ill\_formed},
\texttt{liar\_ill\_formed}; капстоун \texttt{err\_well\_formedness\_synthesis}.\\
\emph{Файл:} \texttt{ERRWellFormedness.v}.
\end{itemize}
\normalsize

% --- Дальше: Блок 3 (КАТЕГОРИЯ) по тезисному плану в Книги/Теория Систем/. ---

\end{document}
